\begin{abstract}
    Retrieval-Augmented Generation (RAG) is an effective approach for legal question answering, yet standard RAG struggles with the lexical gap between colloquial fact patterns and formal statutory corpora. While retrieving with pseudo-documents (e.g., HyDE) proves successful in general-domain RAG, it fails in legal contexts due to this lexical mismatch. To address this, we introduce \hyre{} (Snap-answer COnditioned Pseudo-document Embedding), which first generates a direct "snap answer" reflecting immediate legal intuition. A pseudo-document is then generated in the style of formal legal authority to support this snap answer. Crucially, only the pseudo-document is embedded for retrieval. During the final generation phase, the snap answer and pseudo-document are discarded, ensuring the ultimate response is grounded strictly in the retrieved evidence. We evaluate \hyre{} across two extremes of the question-corpus lexical gap. On BarExamQA, with low lexical similarity, \hyre{} significantly bridges this gap, lifting gold passage retrieval (Hit@5) from a 1.4\% baseline to 9.5\%–12.1\% across 8B, 26B, and 70B parameter models, while consistently improving final answer accuracy. These results demonstrate that \hyre{} is a robust and tailored solution for bridging the colloquial-to-statutory vocabulary gap in complex Legal RAG challenges.
    \end{abstract}